% from ega0-19-fr.tex
\oldpage[0\textsubscript{IV-1}]{79}

\begin{proof}
On vient de voir que la condition est suffisante. Inversement, supposons $M$
formellement projectif et soit $(\mathfrak J_\lambda)$ un système fondamental
d'idéaux ouverts dans $A$. Pour tout $\lambda$, soient $P$ un
$(A/\mathfrak J_\lambda)$-module libre et
$u:P\to M/\mathfrak J_\lambda M$ un
$(A/\mathfrak J_\lambda)$-homomorphisme surjectif. Il existe donc un
$\mathfrak J_\mu\subset\mathfrak J_\lambda$ tel que l'homomorphisme canonique
$M/\mathfrak J_\mu M\to M/\mathfrak J_\lambda M$ se factorise en
\[
  M/\mathfrak J_\mu M\xrightarrow{w}P\xrightarrow{u}M/\mathfrak J_\lambda M;
\]
mais comme
$w(\mathfrak J_\lambda(M/\mathfrak J_\mu M))\subset\mathfrak J_\lambda P=0$,
$w$ se factorise en $M/\mathfrak J_\lambda M\xrightarrow{v}P$, et il est
clair que $u\circ v$ est l'identité dans $M/\mathfrak J_\lambda M$, ce qui
prouve que ce $(A/\mathfrak J_\lambda)$-module est projectif.
\end{proof}

\begin{proposition}[19.2.5]
\label{0.19.2.5-fr}
Soient $A$ un anneau topologique, $M$ un $A$-module topologique.
\begin{enumerate}
  \item[(i)] Pour que $M$ soit formellement projectif (resp. strictement
  formellement projectif), il faut et il suffit que le $\hat A$-module
  topologique $\hat M$ le soit.
  \item[(ii)] Soit $A'$ une $A$-algèbre topologique. Si $M$ est formellement
  projectif (resp. strictement formellement projectif), alors
  $M\otimes_A A'$ (muni de la topologie produit tensoriel) est un $A'$-module
  topologique formellement projectif (resp. strictement formellement
  projectif).
\end{enumerate}
\end{proposition}

\begin{proof}
\begin{enumerate}
  \item[(i)] Il suffit de remarquer que lorsque $\mathfrak J$ (resp. $V$)
  parcourt l'ensemble des idéaux ouverts de $A$ (resp. l'ensemble des
  sous-modules ouverts de $M$), le séparé complété $\hat{\mathfrak J}$
  (resp. $\hat V$) parcourt l'ensemble des idéaux ouverts de $\hat A$
  (resp. l'ensemble des sous-modules ouverts de $\hat M$), et
  $\hat A/\hat{\mathfrak J}=A/\mathfrak J$ (resp.
  $\hat M/\hat V=M/V$) à un isomorphisme canonique près~; comme la notion
  de module formellement projectif (resp. strictement formellement projectif)
  ne fait intervenir que les $A/\mathfrak J$ et les $M/V$, on en déduit
  aussitôt (i).
  \item[(ii)] Supposons d'abord $M$ formellement projectif et posons
  $M'=M\otimes_A A'$~; soient $\mathfrak J'$ un idéal ouvert de $A'$, $P'$,
  $Q'$ deux $(A'/\mathfrak J')$-modules discrets, $u':P'\to Q'$ un
  $A'$-homomorphisme surjectif, $v':M'\to Q'$ un $A'$-homomorphisme continu.
  Il y a un idéal ouvert $\mathfrak J$ de $A$ tel que
  $\mathfrak JA'\subset\mathfrak J'$, donc $P'$ et $Q'$ sont aussi des
  $(A/\mathfrak J)$-modules discrets. Si l'on considère le
  $A$-homomorphisme composé $v:M\xrightarrow{j}M'\xrightarrow{v'}Q'$ qui est
  continu, l'hypothèse entraîne qu'il existe un $A$-homomorphisme continu
  $w:M\to P'$ tel que $v=u'\circ w$~; mais comme $P'$ est un $A'$-module
  topologique, $w$ se factorise en
  $M\xrightarrow{j}M'\xrightarrow{w'}P'$, où $w'$ est un $A'$-homomorphisme
  continu, et comme $v'\circ j=(u'\circ w')\circ j$, on en conclut que
  $v'=u'\circ w'$.

  Supposons ensuite que $M$ soit strictement formellement projectif~; soit
  $(\mathfrak J_\lambda)$ (resp. $(V_\lambda)$) un système fondamental
  d'idéaux ouverts (resp. de sous-modules ouverts) dans $A$ (resp. $M$) tel
  que $(M/V_\lambda)$ soit un $(A/\mathfrak J_\lambda)$-module projectif,
  et soit $(\mathfrak J'_\mu)$ un système fondamental d'idéaux ouverts dans
  $A'$. On a un système fondamental de voisinages de $0$ dans $M'$ en prenant
  les sous-modules
  \[
    \operatorname{Im}(M\otimes_A\mathfrak J'_\mu)
    +\operatorname{Im}(V_\lambda\otimes_A A')=W_{\lambda\mu},
  \]
  où $\lambda$ et $\mu$ sont tels que
  $\mathfrak J_\lambda A'\subset\mathfrak J'_\mu$. Comme
  \[
    (M\otimes_A A')/W_{\lambda\mu}
    =(M/V_\lambda)\otimes_{A/\mathfrak J_\lambda}
      (A'/\mathfrak J'_\mu)
  \]
  et que $(M/V_\lambda)$ est un
  $(A/\mathfrak J_\lambda)$-module projectif, $M'/W_{\lambda\mu}$
  est un $(A'/\mathfrak J'_\mu)$-module projectif.
\end{enumerate}
\end{proof}

\subsection{Algèbres formellement lisses}
\label{subsection:0.19.3-fr}

\begin{definition}[19.3.1]
\label{0.19.3.1-fr}
Soient $A$ un anneau topologique, $B$ une $A$-algèbre topologique. On dit que
$B$ est une $A$-algèbre \emph{formellement lisse} si, pour toute
$A$-algèbre topologique discrète $C$ et tout idéal nilpotent $\mathfrak J$ de
$C$, tout $A$-homomorphisme continu $u:B\to C/\mathfrak J$ se factorise en
\[
  B\xrightarrow{v}C\xrightarrow{\varphi}C/\mathfrak J,
\]
où $v$ est un homomorphisme continu et $\varphi$ l'homomorphisme canonique.
\end{definition}

